\section{Introduction}
\label{sec:intro}

Online job advertisements describe the knowledge and skills employers seek and provide evidence of changing occupational requirements~\citep{ILO_2020_OJVs_BigData,senger2024survey}. Document-level methods retrieve skills from predefined inventories~\citep{bhola2020retrieving}, whereas span extraction locates requirement mentions in the source text~\citep{gnehm2022skills}. Job Skill Named Entity Recognition (JobSkillNER) combines span identification with type assignment. Preserving source wording allows readers to inspect span boundaries and types in context.

Even short Chinese phrases can require multiple decisions about span boundaries and types. A technology name may refer to tool use in one clause and knowledge of underlying principles in another. Coordinated expressions may describe separate activities or share an action that must be retained with its objects. Chinese text generally lacks spaces between words, and tokenization alone cannot resolve these semantic boundaries~\citep{peng2015weibo,chen2021boundary}. Although similar ambiguities occur in other languages~\citep{nguyen2024rethinking}, translating an existing dataset does not by itself determine which Chinese substrings should be annotated. Explicit, context-sensitive annotation rules are therefore needed.

Existing resources provide useful foundations, but their annotation schemes and prediction targets differ. English SkillSpan permits overlapping SKILL and KNOWLEDGE annotation layers~\citep{zhang2022skillspan}. Kompetencer adds fine-grained labels derived from the European Skills, Competences, Qualifications and Occupations (ESCO) classification to annotated Danish and English spans through distant supervision~\citep{jensen2022kompetencer}. Work on German job advertisements extracts education, experience, and language requirements and maps their skill content to taxonomy concepts~\citep{gnehm2022skills}. Document-level skill retrieval and Chinese skill-demand forecasting address different prediction targets~\citep{bhola2020retrieving,chen2024jobsdf}. These differences motivate a Chinese resource with explicit conventions for span boundaries and competency types, coupled with a shared protocol for comparing extraction systems.

We introduce Chinese-SkillSpan, a benchmark for competency span extraction in Chinese job advertisements. It brings together texts from multiple recruitment collections, character-level annotations, an operational handbook, and a shared evaluation protocol (Figure~\ref{fig:macro-micro}). The handbook distinguishes language (L), knowledge-related requirements (K), occupational skills (S), and transversal competencies (T). We use \emph{competency spans} for contiguous source-text mentions across these categories and \emph{occupational skills} for category S. Human coding and adjudication informed the handbook and reference annotations, while large language models generated Silver labels for supervision.

The evaluation asks how consistently annotators apply the handbook and how extraction models perform under its conventions. A separate blinded annotation study assesses agreement under the final handbook, and sampled human review examines Silver-label quality. Model experiments compare prompted language models, Qwen with and without Silver-supervised adaptation, and multilingual encoders trained with Chinese supervision. Diagnostic analyses examine competency types, span lengths, and false positives in sentences with no reference spans. The human reference also informed handbook development, so these model comparisons serve a diagnostic purpose.

The study contributes a Chinese source-text resource with distinct human-reference and Silver annotation layers, an operational handbook for Chinese boundary and type decisions, and evidence on annotation consistency and Silver-label quality. The handbook makes context-sensitive decisions explicit through rules and positive and negative examples. Shared scoring procedures support diagnostic model comparisons and analysis of boundary and type errors.
